\documentclass[11pt]{article}
\usepackage{amsmath, amssymb, geometry, graphicx, xcolor}
\geometry{letterpaper, margin=1in}

\title{Homework 4.9: Free Mechanical Vibrations}
\author{Tashfeen Omran}
\date{September 2026}

\begin{document}

\maketitle
\tableofcontents
\newpage

\section{Part 1: Comprehensive Introduction, Context, and Prerequisites}

\subsection{Core Concepts}
The topic of \textbf{Free Mechanical Vibrations} deals with modeling the motion of a mass attached to a spring. When a mass is displaced from its equilibrium position and released, it oscillates. This system is governed by three primary forces:
\begin{itemize}
    \item \textbf{Restoring Force (Hooke's Law):} The spring exerts a force proportional to the displacement $y$, acting in the opposite direction: $F_s = -ky$, where $k > 0$ is the spring stiffness constant.
    \item \textbf{Damping Force:} Frictional forces (like air resistance or a dashpot) oppose the motion. This force is typically proportional to the velocity: $F_d = -by'$, where $b \ge 0$ is the damping constant.
    \item \textbf{Inertial Force:} According to Newton's Second Law, the net force equals mass times acceleration: $F_{\text{net}} = ma = my''$.
\end{itemize}
Equating the forces yields the fundamental second-order linear homogeneous differential equation for free vibrations:
\[ my'' + by' + ky = 0 \]

\subsection{Intuition and Derivation}
The equation $my'' + by' + ky = 0$ perfectly encapsulates the balance of energy in a mechanical system. If $b = 0$ (undamped), the system continuously trades potential energy (stored in the spring) for kinetic energy (movement of the mass), resulting in perpetual sinusoidal motion. When $b > 0$ (damped), energy is dissipated as heat, causing the oscillations to decay over time. The nature of this decay depends heavily on the discriminant of the characteristic equation, $b^2 - 4mk$.

\subsection{Historical Context and Motivation}
The mathematical modeling of vibrations began in the 17th century. Robert Hooke published his law of elasticity in 1660, and Isaac Newton formulated his laws of motion in 1687. Mathematicians like Leonhard Euler and Daniel Bernoulli later applied calculus to these physical laws to solve differential equations. The motivation was highly practical: understanding vibrations was essential for designing accurate pendulum clocks, building stable bridges, and understanding acoustics and musical instruments.

\subsection{Key Formulas}
\begin{itemize}
    \item \textbf{General Equation:} $my'' + by' + ky = 0$
    \item \textbf{Undamped Motion ($b = 0$):} 
    \[ y(t) = C_1 \cos(\omega t) + C_2 \sin(\omega t) = A \sin(\omega t + \phi) \]
    where $\omega = \sqrt{k/m}$ is the \textbf{angular frequency}.
    \item \textbf{Amplitude and Phase Angle:} 
    \[ A = \sqrt{C_1^2 + C_2^2}, \quad \tan \phi = \frac{C_1}{C_2} \]
    \textit{Note: The quadrant of $\phi$ is determined by the signs of $C_1$ (which corresponds to $\sin \phi$) and $C_2$ (which corresponds to $\cos \phi$).}
    \item \textbf{Period and Natural Frequency:} Period $P = \frac{2\pi}{\omega}$, Frequency $f = \frac{\omega}{2\pi}$.
    \item \textbf{Damped Motion ($b > 0$):} Roots of the auxiliary equation $mr^2 + br + k = 0$ dictate the behavior:
    \begin{itemize}
        \item \textbf{Underdamped ($b^2 < 4mk$):} Complex roots $\alpha \pm i\beta$. $y(t) = A e^{\alpha t} \sin(\beta t + \phi)$.
        \item \textbf{Critically Damped ($b^2 = 4mk$):} Repeated real root $r = -b/2m$. $y(t) = (C_1 + C_2 t)e^{rt}$.
        \item \textbf{Overdamped ($b^2 > 4mk$):} Distinct real roots $r_1, r_2$. $y(t) = C_1 e^{r_1 t} + C_2 e^{r_2 t}$.
    \end{itemize}
\end{itemize}

\subsection{Prerequisites}
To master this topic, you must be proficient in:
\begin{itemize}
    \item Solving second-order linear homogeneous differential equations with constant coefficients.
    \item Basic trigonometry, specifically the addition formulas and the arctangent function.
    \item Differential calculus (product rule, chain rule) for finding maximum displacements (where $y'(t) = 0$).
    \item Algebraic manipulation of complex numbers and radicals.
\end{itemize}

\newpage
\section{Part 2: Detailed Homework Solutions}

\subsection{Problem 1}
\textbf{Statement:} A $2\text{-kg}$ mass is attached to a spring with stiffness $k = 50 \text{ N/m}$. The mass is displaced $1/4 \text{ m}$ to the left of the equilibrium point and given a velocity of $1 \text{ m/sec}$ to the left. Neglecting damping, find the equation of motion of the mass along with the amplitude, period, and frequency. How long after release does the mass pass through the equilibrium position?

\textbf{Solution:}
\begin{enumerate}
    \item \textbf{Set up the Initial Value Problem (IVP):}
    We have $m = 2$, $k = 50$, and $b = 0$. The initial displacement is to the left, so $y(0) = -1/4$. The initial velocity is to the left, so $y'(0) = -1$.
    \[ 2y'' + 50y = 0 \implies y'' + 25y = 0 \]
    
    \item \textbf{Solve the Differential Equation:}
    The characteristic equation is $r^2 + 25 = 0$, which gives roots $r = \pm 5i$.
    The general solution is:
    \[ y(t) = C_1 \cos(5t) + C_2 \sin(5t) \]
    
    \item \textbf{Apply Initial Conditions:}
    \[ y(0) = C_1 \cos(0) + C_2 \sin(0) = C_1 = -\frac{1}{4} \]
    Take the derivative:
    \[ y'(t) = -5C_1 \sin(5t) + 5C_2 \cos(5t) \]
    \[ y'(0) = 5C_2 = -1 \implies C_2 = -\frac{1}{5} \]
    The equation of motion is:
    \[ y(t) = -0.25 \cos(5t) - 0.2 \sin(5t) \]
    
    \item \textbf{Find Amplitude, Period, and Frequency:}
    Amplitude $A = \sqrt{C_1^2 + C_2^2} = \sqrt{\left(-\frac{1}{4}\right)^2 + \left(-\frac{1}{5}\right)^2} = \sqrt{\frac{1}{16} + \frac{1}{25}} = \sqrt{\frac{41}{400}} = \frac{\sqrt{41}}{20}$.
    Period $P = \frac{2\pi}{\omega} = \frac{2\pi}{5} \text{ seconds}$.
    Frequency $f = \frac{\omega}{2\pi} = \frac{5}{2\pi} \text{ cycles/sec}$.
    
    \item \textbf{Convert to Phase-Amplitude Form to find Equilibrium Time:}
    We want $y(t) = A \sin(5t + \phi)$.
    $\tan \phi = \frac{C_1}{C_2} = \frac{-1/4}{-1/5} = \frac{5}{4} = 1.25$.
    Since $C_1 < 0$ ($\sin \phi < 0$) and $C_2 < 0$ ($\cos \phi < 0$), $\phi$ is in Quadrant III.
    \[ \phi = \arctan(1.25) + \pi \approx 0.896 + 3.142 = 4.038 \text{ radians} \]
    The equation is $y(t) = \frac{\sqrt{41}}{20} \sin(5t + 4.038)$.
    To find when it passes through equilibrium, set $y(t) = 0$:
    \[ \sin(5t + 4.038) = 0 \implies 5t + 4.038 = n\pi \]
    For the first positive time, we choose $n = 2$ (since $n=1$ gives a negative time):
    \[ 5t + 4.038 = 2\pi \approx 6.283 \implies 5t = 2.245 \implies t \approx 0.449 \text{ seconds} \]
\end{enumerate}
\textbf{Final Answer:} $y(t) = -0.25 \cos(5t) - 0.2 \sin(5t)$. $A = \frac{\sqrt{41}}{20}$, $P = \frac{2\pi}{5}$, $f = \frac{5}{2\pi}$. Time to equilibrium is $t \approx 0.449 \text{ s}$.

\subsection{Problem 2}
\textbf{Statement:} A $3\text{-kg}$ mass is attached to a spring with stiffness $k = 48 \text{ N/m}$. The mass is displaced $1/2 \text{ m}$ to the left of the equilibrium point and given a velocity of $2 \text{ m/sec}$ to the right. Find the equation of motion, amplitude, period, and frequency. How long after release does the mass pass through the equilibrium position?

\textbf{Solution:}
\begin{enumerate}
    \item \textbf{Set up the IVP:}
    $m = 3$, $k = 48$, $b = 0$. $y(0) = -1/2$ (left), $y'(0) = 2$ (right).
    \[ 3y'' + 48y = 0 \implies y'' + 16y = 0 \]
    
    \item \textbf{Solve the Differential Equation:}
    Characteristic equation: $r^2 + 16 = 0 \implies r = \pm 4i$.
    \[ y(t) = C_1 \cos(4t) + C_2 \sin(4t) \]
    
    \item \textbf{Apply Initial Conditions:}
    \[ y(0) = C_1 = -\frac{1}{2} \]
    \[ y'(t) = -4C_1 \sin(4t) + 4C_2 \cos(4t) \implies y'(0) = 4C_2 = 2 \implies C_2 = \frac{1}{2} \]
    Equation of motion: $y(t) = -0.5 \cos(4t) + 0.5 \sin(4t)$.
    
    \item \textbf{Find Amplitude, Period, and Frequency:}
    $A = \sqrt{(-1/2)^2 + (1/2)^2} = \sqrt{1/4 + 1/4} = \sqrt{1/2} = \frac{\sqrt{2}}{2}$.
    Period $P = \frac{2\pi}{4} = \frac{\pi}{2} \text{ seconds}$.
    Frequency $f = \frac{4}{2\pi} = \frac{2}{\pi} \text{ cycles/sec}$.
    
    \item \textbf{Convert to Phase-Amplitude Form:}
    $\tan \phi = \frac{C_1}{C_2} = \frac{-1/2}{1/2} = -1$.
    Since $C_1 < 0$ ($\sin \phi < 0$) and $C_2 > 0$ ($\cos \phi > 0$), $\phi$ is in Quadrant IV.
    \[ \phi = -\frac{\pi}{4} \]
    \[ y(t) = \frac{\sqrt{2}}{2} \sin\left(4t - \frac{\pi}{4}\right) \]
    Set $y(t) = 0$ for equilibrium:
    \[ 4t - \frac{\pi}{4} = 0 \implies 4t = \frac{\pi}{4} \implies t = \frac{\pi}{16} \text{ seconds} \]
\end{enumerate}
\textbf{Final Answer:} $y(t) = -0.5 \cos(4t) + 0.5 \sin(4t)$. $A = \frac{\sqrt{2}}{2}$, $P = \frac{\pi}{2}$, $f = \frac{2}{\pi}$. Time to equilibrium is $t = \frac{\pi}{16} \approx 0.196 \text{ s}$.

\subsection{Problem 3}
\textbf{Statement:} The motion of a mass-spring system with damping is governed by $y''(t) + by'(t) + 16y(t) = 0$; $y(0) = 1$, $y'(0) = 0$. Find the equation of motion for $b = 0, 6, 8, \text{ and } 10$.

\textbf{Solution:}
\begin{enumerate}
    \item \textbf{Case 1: $b = 0$ (Undamped)}
    \[ y'' + 16y = 0 \implies r^2 + 16 = 0 \implies r = \pm 4i \]
    \[ y(t) = C_1 \cos(4t) + C_2 \sin(4t) \]
    $y(0) = C_1 = 1$. $y'(0) = 4C_2 = 0 \implies C_2 = 0$.
    \textbf{Answer:} $y(t) = \cos(4t)$.
    
    \item \textbf{Case 2: $b = 6$ (Underdamped)}
    \[ y'' + 6y' + 16y = 0 \implies r^2 + 6r + 16 = 0 \]
    \[ r = \frac{-6 \pm \sqrt{36 - 64}}{2} = -3 \pm i\sqrt{7} \]
    \[ y(t) = e^{-3t}(C_1 \cos(\sqrt{7}t) + C_2 \sin(\sqrt{7}t)) \]
    $y(0) = C_1 = 1$.
    $y'(t) = -3e^{-3t}(C_1 \cos(\sqrt{7}t) + C_2 \sin(\sqrt{7}t)) + e^{-3t}(-\sqrt{7}C_1 \sin(\sqrt{7}t) + \sqrt{7}C_2 \cos(\sqrt{7}t))$.
    $y'(0) = -3C_1 + \sqrt{7}C_2 = 0 \implies -3(1) + \sqrt{7}C_2 = 0 \implies C_2 = \frac{3}{\sqrt{7}}$.
    \textbf{Answer:} $y(t) = e^{-3t}\left(\cos(\sqrt{7}t) + \frac{3}{\sqrt{7}}\sin(\sqrt{7}t)\right)$.
    
    \item \textbf{Case 3: $b = 8$ (Critically Damped)}
    \[ y'' + 8y' + 16y = 0 \implies r^2 + 8r + 16 = 0 \implies (r+4)^2 = 0 \implies r = -4 \]
    \[ y(t) = (C_1 + C_2 t)e^{-4t} \]
    $y(0) = C_1 = 1$.
    $y'(t) = C_2 e^{-4t} - 4(C_1 + C_2 t)e^{-4t}$.
    $y'(0) = C_2 - 4C_1 = 0 \implies C_2 = 4(1) = 4$.
    \textbf{Answer:} $y(t) = (1 + 4t)e^{-4t}$.
    
    \item \textbf{Case 4: $b = 10$ (Overdamped)}
    \[ y'' + 10y' + 16y = 0 \implies r^2 + 10r + 16 = 0 \implies (r+2)(r+8) = 0 \implies r = -2, -8 \]
    \[ y(t) = C_1 e^{-2t} + C_2 e^{-8t} \]
    $y(0) = C_1 + C_2 = 1$.
    $y'(t) = -2C_1 e^{-2t} - 8C_2 e^{-8t} \implies y'(0) = -2C_1 - 8C_2 = 0 \implies C_1 = -4C_2$.
    Substitute into first eq: $-4C_2 + C_2 = 1 \implies -3C_2 = 1 \implies C_2 = -1/3$.
    $C_1 = -4(-1/3) = 4/3$.
    \textbf{Answer:} $y(t) = \frac{4}{3}e^{-2t} - \frac{1}{3}e^{-8t}$.
\end{enumerate}

\subsection{Problem 6}
\textbf{Statement:} The motion of a mass-spring system with damping is governed by $y''(t) + 4y'(t) + ky(t) = 0$; $y(0) = 1$, $y'(0) = 0$. Find the equation of motion for $k = 2, 4, \text{ and } 6$.

\textbf{Solution:}
\begin{enumerate}
    \item \textbf{Case 1: $k = 2$ (Overdamped)}
    \[ r^2 + 4r + 2 = 0 \implies r = \frac{-4 \pm \sqrt{16 - 8}}{2} = -2 \pm \sqrt{2} \]
    \[ y(t) = C_1 e^{(-2+\sqrt{2})t} + C_2 e^{(-2-\sqrt{2})t} \]
    $y(0) = C_1 + C_2 = 1$.
    $y'(0) = (-2+\sqrt{2})C_1 + (-2-\sqrt{2})C_2 = 0$.
    Rearranging: $-2(C_1 + C_2) + \sqrt{2}(C_1 - C_2) = 0 \implies -2(1) + \sqrt{2}(C_1 - C_2) = 0 \implies C_1 - C_2 = \sqrt{2}$.
    Adding the two equations: $2C_1 = 1 + \sqrt{2} \implies C_1 = \frac{1+\sqrt{2}}{2}$.
    Subtracting: $2C_2 = 1 - \sqrt{2} \implies C_2 = \frac{1-\sqrt{2}}{2}$.
    \textbf{Answer:} $y(t) = \frac{1+\sqrt{2}}{2} e^{(-2+\sqrt{2})t} + \frac{1-\sqrt{2}}{2} e^{(-2-\sqrt{2})t}$.
    
    \item \textbf{Case 2: $k = 4$ (Critically Damped)}
    \[ r^2 + 4r + 4 = 0 \implies (r+2)^2 = 0 \implies r = -2 \]
    \[ y(t) = (C_1 + C_2 t)e^{-2t} \]
    $y(0) = C_1 = 1$.
    $y'(0) = C_2 - 2C_1 = 0 \implies C_2 = 2$.
    \textbf{Answer:} $y(t) = (1 + 2t)e^{-2t}$.
    
    \item \textbf{Case 3: $k = 6$ (Underdamped)}
    \[ r^2 + 4r + 6 = 0 \implies r = \frac{-4 \pm \sqrt{16 - 24}}{2} = -2 \pm i\sqrt{2} \]
    \[ y(t) = e^{-2t}(C_1 \cos(\sqrt{2}t) + C_2 \sin(\sqrt{2}t)) \]
    $y(0) = C_1 = 1$.
    $y'(0) = -2C_1 + \sqrt{2}C_2 = 0 \implies -2(1) + \sqrt{2}C_2 = 0 \implies C_2 = \frac{2}{\sqrt{2}} = \sqrt{2}$.
    \textbf{Answer:} $y(t) = e^{-2t}(\cos(\sqrt{2}t) + \sqrt{2}\sin(\sqrt{2}t))$.
\end{enumerate}

\subsection{Problem 7}
\textbf{Statement:} A $1/8\text{-kg}$ mass is attached to a spring with stiffness $16 \text{ N/m}$. The damping constant is $2 \text{ N-sec/m}$. The mass is moved $3/4 \text{ m}$ to the left and given an initial velocity of $2 \text{ m/sec}$ to the right. Determine the equation of motion, damping factor, quasiperiod, and quasifrequency.

\textbf{Solution:}
\begin{enumerate}
    \item \textbf{Set up the IVP:}
    $m = 1/8$, $b = 2$, $k = 16$. $y(0) = -3/4$, $y'(0) = 2$.
    \[ \frac{1}{8}y'' + 2y' + 16y = 0 \implies y'' + 16y' + 128y = 0 \]
    
    \item \textbf{Solve the Differential Equation:}
    \[ r^2 + 16r + 128 = 0 \implies r = \frac{-16 \pm \sqrt{256 - 512}}{2} = -8 \pm 8i \]
    \[ y(t) = e^{-8t}(C_1 \cos(8t) + C_2 \sin(8t)) \]
    
    \item \textbf{Apply Initial Conditions:}
    \[ y(0) = C_1 = -\frac{3}{4} \]
    \[ y'(t) = -8e^{-8t}(C_1 \cos(8t) + C_2 \sin(8t)) + e^{-8t}(-8C_1 \sin(8t) + 8C_2 \cos(8t)) \]
    \[ y'(0) = -8C_1 + 8C_2 = 2 \implies -8\left(-\frac{3}{4}\right) + 8C_2 = 2 \implies 6 + 8C_2 = 2 \implies 8C_2 = -4 \implies C_2 = -\frac{1}{2} \]
    Equation of motion: $y(t) = e^{-8t}\left(-\frac{3}{4} \cos(8t) - \frac{1}{2} \sin(8t)\right)$.
    
    \item \textbf{Find Damping Factor, Quasiperiod, and Quasifrequency:}
    The damping factor is $A e^{\alpha t}$.
    $A = \sqrt{C_1^2 + C_2^2} = \sqrt{\left(-\frac{3}{4}\right)^2 + \left(-\frac{1}{2}\right)^2} = \sqrt{\frac{9}{16} + \frac{4}{16}} = \frac{\sqrt{13}}{4}$.
    Damping factor: $\frac{\sqrt{13}}{4} e^{-8t}$.
    Quasiperiod $P = \frac{2\pi}{\beta} = \frac{2\pi}{8} = \frac{\pi}{4} \text{ seconds}$.
    Quasifrequency $f = \frac{\beta}{2\pi} = \frac{8}{2\pi} = \frac{4}{\pi} \text{ cycles/sec}$.
\end{enumerate}
\textbf{Final Answer:} $y(t) = e^{-8t}\left(-\frac{3}{4} \cos(8t) - \frac{1}{2} \sin(8t)\right)$. Damping factor is $\frac{\sqrt{13}}{4} e^{-8t}$, quasiperiod is $\frac{\pi}{4}$, quasifrequency is $\frac{4}{\pi}$.

\subsection{Problem 8}
\textbf{Statement:} A $20\text{-kg}$ mass is attached to a spring with stiffness $200 \text{ N/m}$. The damping constant is $140 \text{ N-sec/m}$. If the mass is pulled $25 \text{ cm}$ to the right and given an initial leftward velocity of $1 \text{ m/sec}$, when will it first return to its equilibrium position?

\textbf{Solution:}
\begin{enumerate}
    \item \textbf{Set up the IVP:}
    $m = 20$, $b = 140$, $k = 200$. $y(0) = 0.25 \text{ m}$ (right), $y'(0) = -1 \text{ m/sec}$ (left).
    \[ 20y'' + 140y' + 200y = 0 \implies y'' + 7y' + 10y = 0 \]
    
    \item \textbf{Solve the Differential Equation:}
    \[ r^2 + 7r + 10 = 0 \implies (r+2)(r+5) = 0 \implies r = -2, -5 \]
    \[ y(t) = C_1 e^{-2t} + C_2 e^{-5t} \]
    
    \item \textbf{Apply Initial Conditions:}
    \[ y(0) = C_1 + C_2 = 0.25 = \frac{1}{4} \]
    \[ y'(0) = -2C_1 - 5C_2 = -1 \]
    Multiply the first equation by 2: $2C_1 + 2C_2 = 0.5$.
    Add to the second equation: $-3C_2 = -0.5 \implies C_2 = \frac{1}{6}$.
    $C_1 = \frac{1}{4} - \frac{1}{6} = \frac{3}{12} - \frac{2}{12} = \frac{1}{12}$.
    Equation of motion: $y(t) = \frac{1}{12}e^{-2t} + \frac{1}{6}e^{-5t}$.
    
    \item \textbf{Find Return to Equilibrium:}
    Set $y(t) = 0$:
    \[ \frac{1}{12}e^{-2t} + \frac{1}{6}e^{-5t} = 0 \implies e^{-2t} + 2e^{-5t} = 0 \implies e^{3t} = -2 \]
    Since $e^{3t}$ is always positive, there is no real solution for $t$.
\end{enumerate}
\textbf{Final Answer:} The mass will \textbf{never} return to its equilibrium position.

\subsection{Problem 9}
\textbf{Statement:} A $2\text{-kg}$ mass is attached to a spring with stiffness $40 \text{ N/m}$. The damping constant is $8\sqrt{5} \text{ N-sec/m}$. If the mass is pulled $10 \text{ cm}$ to the right and given an initial rightward velocity of $2 \text{ m/sec}$, what is the maximum displacement from equilibrium that it will attain?

\textbf{Solution:}
\begin{enumerate}
    \item \textbf{Set up the IVP:}
    $m = 2$, $b = 8\sqrt{5}$, $k = 40$. $y(0) = 0.1 \text{ m}$, $y'(0) = 2 \text{ m/sec}$.
    \[ 2y'' + 8\sqrt{5}y' + 40y = 0 \implies y'' + 4\sqrt{5}y' + 20y = 0 \]
    
    \item \textbf{Solve the Differential Equation:}
    \[ r^2 + 4\sqrt{5}r + 20 = 0 \]
    Discriminant: $(4\sqrt{5})^2 - 4(1)(20) = 16(5) - 80 = 0$. The system is critically damped.
    \[ r = -2\sqrt{5} \]
    \[ y(t) = (C_1 + C_2 t)e^{-2\sqrt{5}t} \]
    
    \item \textbf{Apply Initial Conditions:}
    \[ y(0) = C_1 = 0.1 \]
    \[ y'(t) = C_2 e^{-2\sqrt{5}t} - 2\sqrt{5}(C_1 + C_2 t)e^{-2\sqrt{5}t} \]
    \[ y'(0) = C_2 - 2\sqrt{5}C_1 = 2 \implies C_2 - 0.2\sqrt{5} = 2 \implies C_2 = 2 + 0.2\sqrt{5} \]
    
    \item \textbf{Find Maximum Displacement:}
    Maximum displacement occurs when $y'(t) = 0$:
    \[ C_2 - 2\sqrt{5}C_1 - 2\sqrt{5}C_2 t = 0 \]
    We know $C_2 - 2\sqrt{5}C_1 = y'(0) = 2$.
    \[ 2 - 2\sqrt{5}(2 + 0.2\sqrt{5})t = 0 \implies 2 - (4\sqrt{5} + 2)t = 0 \implies t = \frac{2}{4\sqrt{5} + 2} = \frac{1}{2\sqrt{5} + 1} \]
    Rationalizing the denominator:
    \[ t = \frac{2\sqrt{5} - 1}{(2\sqrt{5})^2 - 1^2} = \frac{2\sqrt{5} - 1}{19} \text{ seconds} \]
    Substitute $t$ back into $y(t)$:
    \[ y_{\text{max}} = \left(0.1 + (2 + 0.2\sqrt{5})\left(\frac{2\sqrt{5} - 1}{19}\right)\right) e^{-2\sqrt{5}\left(\frac{2\sqrt{5} - 1}{19}\right)} \]
    Simplify the coefficient:
    \[ 0.1 + \left(\frac{10 + \sqrt{5}}{5}\right)\left(\frac{2\sqrt{5} - 1}{19}\right) = \frac{1}{10} + \frac{20\sqrt{5} - 10 + 10 - \sqrt{5}}{95} = \frac{1}{10} + \frac{19\sqrt{5}}{95} = \frac{1}{10} + \frac{\sqrt{5}}{5} = \frac{1 + 2\sqrt{5}}{10} \]
    The exponent is:
    \[ -2\sqrt{5}\left(\frac{2\sqrt{5} - 1}{19}\right) = \frac{-20 + 2\sqrt{5}}{19} \]
\end{enumerate}
\textbf{Final Answer:} The maximum displacement is $y_{\text{max}} = \frac{1 + 2\sqrt{5}}{10} e^{\frac{2\sqrt{5} - 20}{19}} \text{ meters}$.

\subsection{Problem 11}
\textbf{Statement:} A $1\text{-kg}$ mass is attached to a spring with stiffness $100 \text{ N/m}$. The damping constant is $0.2 \text{ N-sec/m}$. If the mass is pushed rightward from equilibrium with a velocity of $1 \text{ m/sec}$, when will it attain its maximum displacement to the right?

\textbf{Solution:}
\begin{enumerate}
    \item \textbf{Set up the IVP:}
    $m = 1$, $b = 0.2$, $k = 100$. $y(0) = 0$, $y'(0) = 1$.
    \[ y'' + 0.2y' + 100y = 0 \]
    
    \item \textbf{Solve the Differential Equation:}
    \[ r^2 + 0.2r + 100 = 0 \implies r = \frac{-0.2 \pm \sqrt{0.04 - 400}}{2} = -0.1 \pm \frac{\sqrt{-399.96}}{2} = -0.1 \pm i\sqrt{99.99} \]
    Let $\beta = \sqrt{99.99}$.
    \[ y(t) = e^{-0.1t}(C_1 \cos(\beta t) + C_2 \sin(\beta t)) \]
    
    \item \textbf{Apply Initial Conditions:}
    \[ y(0) = C_1 = 0 \]
    \[ y'(t) = e^{-0.1t}(\beta C_2 \cos(\beta t)) - 0.1e^{-0.1t}(C_2 \sin(\beta t)) \]
    \[ y'(0) = \beta C_2 = 1 \implies C_2 = \frac{1}{\beta} = \frac{1}{\sqrt{99.99}} \]
    Equation of motion: $y(t) = \frac{1}{\sqrt{99.99}} e^{-0.1t} \sin(\sqrt{99.99}t)$.
    
    \item \textbf{Find Time of Maximum Displacement:}
    Set $y'(t) = 0$:
    \[ \beta \cos(\beta t) - 0.1 \sin(\beta t) = 0 \implies \tan(\beta t) = \frac{\beta}{0.1} = 10\beta = 10\sqrt{99.99} = \sqrt{9999} \]
    \[ \beta t = \arctan(\sqrt{9999}) \implies t = \frac{1}{\sqrt{99.99}} \arctan(\sqrt{9999}) \]
\end{enumerate}
\textbf{Final Answer:} The mass attains its maximum displacement at $t = \frac{1}{\sqrt{99.99}} \arctan(\sqrt{9999}) \approx 0.156 \text{ seconds}$.

\subsection{Problem 14}
\textbf{Statement:} For an underdamped system, verify that as $b \to 0$ the damping factor approaches the constant $A$ and the quasifrequency approaches the natural frequency $\sqrt{k/m}/(2\pi)$.

\textbf{Solution:}
\begin{enumerate}
    \item \textbf{Damping Factor Limit:}
    The damping factor for an underdamped system is given by $A e^{\alpha t}$, where $\alpha = -\frac{b}{2m}$.
    \[ \lim_{b \to 0} A e^{-\frac{b}{2m}t} = A e^{0} = A \]
    Thus, the damping factor approaches the constant amplitude $A$ of an undamped system.
    
    \item \textbf{Quasifrequency Limit:}
    The quasifrequency is $f_q = \frac{\beta}{2\pi}$, where $\beta = \frac{\sqrt{4mk - b^2}}{2m}$.
    \[ \lim_{b \to 0} \frac{\sqrt{4mk - b^2}}{4\pi m} = \frac{\sqrt{4mk}}{4\pi m} = \frac{2\sqrt{mk}}{4\pi m} = \frac{\sqrt{mk}}{2\pi m} = \frac{1}{2\pi} \sqrt{\frac{k}{m}} \]
    This is exactly the natural frequency of the undamped system.
\end{enumerate}
\textbf{Final Answer:} The limits evaluate exactly to $A$ and $\frac{1}{2\pi}\sqrt{\frac{k}{m}}$, verifying the statement.

\subsection{Problem 17}
\textbf{Statement:} Consider the equation for free mechanical vibration, $my'' + by' + ky = 0$, and assume the motion is critically damped. Let $y(0) = y_0$, $y'(0) = v_0$ and assume $y_0 \neq 0$. Prove that the mass will pass through its equilibrium at exactly one positive time if and only if $\frac{-2my_0}{2mv_0 + by_0} > 0$.

\textbf{Solution:}
\begin{enumerate}
    \item \textbf{Set up the Critically Damped Solution:}
    Since it is critically damped, $b^2 = 4mk$, and the root is $r = -\frac{b}{2m}$.
    \[ y(t) = (C_1 + C_2 t)e^{-\frac{b}{2m}t} \]
    
    \item \textbf{Apply Initial Conditions:}
    \[ y(0) = C_1 = y_0 \]
    \[ y'(t) = C_2 e^{-\frac{b}{2m}t} - \frac{b}{2m}(C_1 + C_2 t)e^{-\frac{b}{2m}t} \]
    \[ y'(0) = C_2 - \frac{b}{2m}y_0 = v_0 \implies C_2 = v_0 + \frac{b}{2m}y_0 = \frac{2mv_0 + by_0}{2m} \]
    The equation of motion is:
    \[ y(t) = \left(y_0 + \frac{2mv_0 + by_0}{2m} t\right) e^{-\frac{b}{2m}t} \]
    
    \item \textbf{Find Equilibrium Time:}
    The mass passes through equilibrium when $y(t) = 0$. Since $e^{-\frac{b}{2m}t} \neq 0$:
    \[ y_0 + \frac{2mv_0 + by_0}{2m} t = 0 \implies \frac{2mv_0 + by_0}{2m} t = -y_0 \implies t = \frac{-2my_0}{2mv_0 + by_0} \]
    
    \item \textbf{Condition for Positive Time:}
    For the mass to pass through equilibrium at a \textit{positive} time, we must have $t > 0$.
    Therefore, exactly one positive time exists if and only if:
    \[ \frac{-2my_0}{2mv_0 + by_0} > 0 \]
\end{enumerate}
\textbf{Final Answer:} The proof is complete by directly solving for $t$ and enforcing the condition $t > 0$.

\newpage
\section{Part 3: In-Depth Analysis of Problems and Techniques}

\subsection{A) Problem Types and General Approach}
The homework problems fall into several distinct archetypes:
\begin{itemize}
    \item \textbf{Undamped Initial Value Problems (Problems 1, 2):} 
    \textit{Strategy:} Solve $my'' + ky = 0$ to get $y(t) = C_1 \cos(\omega t) + C_2 \sin(\omega t)$. Use initial conditions to find $C_1$ and $C_2$. Convert to the phase-amplitude form $A \sin(\omega t + \phi)$ to easily find the amplitude, period, and times when $y(t) = 0$.
    \item \textbf{Damped Initial Value Problems (Problems 3, 6, 7, 8):} 
    \textit{Strategy:} Calculate the discriminant $b^2 - 4mk$ to classify the system (underdamped, critically damped, overdamped). Write the corresponding general solution template. Apply $y(0)$ and $y'(0)$ to solve a $2 \times 2$ system of linear equations for $C_1$ and $C_2$.
    \item \textbf{Optimization / Maximum Displacement (Problems 9, 11):} 
    \textit{Strategy:} Once the specific equation of motion $y(t)$ is found, take the first derivative $y'(t)$ and set it to zero. Solve for $t$. Plug this $t$ back into $y(t)$ to find the maximum displacement.
    \item \textbf{Theoretical Proofs and Limits (Problems 14, 17):} 
    \textit{Strategy:} Keep variables symbolic ($m, b, k, y_0, v_0$). Perform the standard IVP steps algebraically. Use calculus limits ($\lim_{b \to 0}$) or algebraic inequalities to prove the required statements.
\end{itemize}

\subsection{B) Key Algebraic and Calculus Manipulations}
\begin{itemize}
    \item \textbf{Phase-Amplitude Conversion (Trigonometric Identity):} Used heavily in Problems 1 and 2. The identity $C_1 \cos(\omega t) + C_2 \sin(\omega t) = A \sin(\omega t + \phi)$ requires calculating $A = \sqrt{C_1^2 + C_2^2}$ and $\tan \phi = C_1/C_2$. \textit{Crucial Trick:} You must check the signs of $C_1$ and $C_2$ to place $\phi$ in the correct quadrant.
    \item \textbf{Product Rule with Chain Rule:} Essential for finding $y'(t)$ in damped systems. For example, in Problem 7, differentiating $e^{-8t}(C_1 \cos 8t + C_2 \sin 8t)$ requires applying the product rule $(uv)' = u'v + uv'$, where the chain rule is used on both the exponential and trigonometric terms.
    \item \textbf{Rationalizing Denominators:} In Problem 9, solving for $t$ yielded $t = \frac{1}{2\sqrt{5} + 1}$. To simplify the final answer, we multiplied the numerator and denominator by the conjugate $2\sqrt{5} - 1$.
    \item \textbf{Solving Exponential Equations:} In Problem 8, setting $y(t) = 0$ led to $e^{3t} = -2$. Recognizing that the range of the real exponential function is $(0, \infty)$ immediately proves that no solution exists, meaning the mass never crosses equilibrium.
\end{itemize}

\newpage
\section{Part 4: "Cheatsheet" and Tips for Success}

\subsection{Summary of Important Formulas}
\begin{itemize}
    \item \textbf{Characteristic Equation:} $mr^2 + br + k = 0 \implies r = \frac{-b \pm \sqrt{b^2 - 4mk}}{2m}$
    \item \textbf{Undamped ($b=0$):} $y(t) = A \sin(\omega t + \phi)$, where $\omega = \sqrt{k/m}$.
    \item \textbf{Underdamped ($b^2 < 4mk$):} $y(t) = A e^{\alpha t} \sin(\beta t + \phi)$, where $\alpha = -b/2m$ and $\beta = \frac{\sqrt{4mk - b^2}}{2m}$.
    \item \textbf{Critically Damped ($b^2 = 4mk$):} $y(t) = (C_1 + C_2 t)e^{-bt/2m}$.
    \item \textbf{Overdamped ($b^2 > 4mk$):} $y(t) = C_1 e^{r_1 t} + C_2 e^{r_2 t}$.
\end{itemize}

\subsection{Tricks and Shortcuts}
\begin{itemize}
    \item \textbf{The Quadrant Cheat Code for $\phi$:} When converting to $A \sin(\omega t + \phi)$, remember that $C_1 = A \sin \phi$ and $C_2 = A \cos \phi$. 
    \begin{itemize}
        \item If $C_1 > 0, C_2 > 0 \implies$ Quadrant I.
        \item If $C_1 > 0, C_2 < 0 \implies$ Quadrant II ($\phi = \arctan(C_1/C_2) + \pi$).
        \item If $C_1 < 0, C_2 < 0 \implies$ Quadrant III ($\phi = \arctan(C_1/C_2) + \pi$).
        \item If $C_1 < 0, C_2 > 0 \implies$ Quadrant IV ($\phi = \arctan(C_1/C_2)$ or add $2\pi$).
    \end{itemize}
    \item \textbf{Quick Check for Critical Damping:} Before using the quadratic formula, quickly check if $b^2 = 4mk$. If it does, you can immediately jump to the repeated root template, saving time.
\end{itemize}

\subsection{Common Pitfalls to Avoid}
\begin{itemize}
    \item \textbf{Forgetting the $t$ in Critical Damping:} A very common mistake is writing $y(t) = C_1 e^{rt} + C_2 e^{rt}$ instead of $y(t) = (C_1 + C_2 t)e^{rt}$. The former collapses into a single constant $(C_1+C_2)e^{rt}$ and cannot satisfy two initial conditions.
    \item \textbf{Sign Errors in Initial Conditions:} "Displaced to the left" usually means $y(0) < 0$, and "velocity to the right" means $y'(0) > 0$. Always double-check your coordinate system definition.
    \item \textbf{Confusing Frequency and Angular Frequency:} $\omega$ (or $\beta$) is the \textit{angular} frequency (radians/sec). The actual frequency $f$ (cycles/sec or Hertz) is $\omega / 2\pi$. Read the question carefully to see which one is requested.
\end{itemize}

\newpage
\section{Part 5: Conceptual Synthesis and The "Big Picture"}

\subsection{A) Thematic Connections}
The core theme of this topic is \textbf{linear operators and eigenvalues}. The differential equation $my'' + by' + ky = 0$ can be viewed as a linear operator $L = m\frac{d^2}{dt^2} + b\frac{d}{dt} + k$ acting on a function $y(t)$. We are looking for functions in the null space of this operator. By guessing a solution of the form $e^{rt}$, we transform the calculus problem into an algebraic eigenvalue problem: finding the roots of the characteristic polynomial. This theme of transforming differential equations into algebraic equations via exponential functions is one of the most powerful recurring motifs in applied mathematics.

\subsection{B) Forward and Backward Links}
\begin{itemize}
    \item \textbf{Backward Links:} This topic is the direct physical application of the purely mathematical techniques learned in the previous chapter on solving second-order linear homogeneous ODEs. It gives physical meaning to the abstract concepts of real, repeated, and complex roots.
    \item \textbf{Forward Links:} Understanding free vibrations is the absolute bedrock for the next section: \textbf{Forced Vibrations} ($my'' + by' + ky = F(t)$). When an external force is applied, the system's natural frequency interacts with the forcing frequency. If they match, it leads to \textbf{resonance}, a phenomenon critical in engineering (e.g., preventing bridges from collapsing). Furthermore, this exact mathematical framework is used in RLC electrical circuits, where mass becomes inductance, damping becomes resistance, and spring stiffness becomes inverse capacitance.
\end{itemize}

\newpage
\section{Part 6: Abstract Real-World Modeling (AI, ML, \& Quant Finance)}

\subsection{A) Concrete Scenarios}
\begin{enumerate}
    \item \textbf{Machine Learning (Optimization):} In training deep neural networks, standard Gradient Descent can get stuck in local minima or oscillate wildly in ravines. The \textbf{Heavy Ball Method (Momentum)} introduces a physical analogy: imagine a mass sliding down the loss landscape with friction. The "mass" provides inertia, helping the optimizer barrel through shallow local minima, while the "friction" (damping) ensures it eventually settles at the bottom.
    \item \textbf{Quantitative Finance (Mean-Reverting Processes):} Interest rates and volatility are often modeled using the Ornstein-Uhlenbeck process, a stochastic differential equation. This process acts like a stochastic damped spring: it has a "restoring force" pulling the rate back to a long-term mean, and "damping" that prevents it from oscillating infinitely, all while being bombarded by random Brownian motion (analogous to a random external forcing function).
\end{enumerate}

\subsection{B) Model Problem Setup: Momentum in Gradient Descent}
Let us rigorously model the Heavy Ball method as a damped mechanical vibration.
Define the state space as $x(t) \in \mathbb{R}^n$, representing the weights of a neural network at continuous time $t$. Let $f(x): \mathbb{R}^n \to \mathbb{R}$ be a twice-differentiable loss function we wish to minimize.

In standard continuous-time gradient descent, the trajectory follows a first-order ODE: $\dot{x} = -\nabla f(x)$.
To introduce momentum, we model the weights as a particle of mass $m > 0$ moving through a viscous medium with damping coefficient $b > 0$, subject to a conservative force field generated by the potential energy $f(x)$. By Newton's Second Law:
\[ m \ddot{x} + b \dot{x} = -\nabla f(x) \]
Assume the loss function is locally a strictly convex quadratic (a common assumption near a local minimum), $f(x) = \frac{1}{2} x^T H x$, where $H$ is the symmetric positive-definite Hessian matrix. The gradient is $\nabla f(x) = Hx$. The system becomes:
\[ m \ddot{x} + b \dot{x} + Hx = 0 \]
This is a multidimensional system of damped harmonic oscillators. By diagonalizing $H = Q \Lambda Q^T$, we can decouple this into $n$ independent 1D equations:
\[ m \ddot{z}_i + b \dot{z}_i + \lambda_i z_i = 0 \quad \text{for } i = 1, \dots, n \]
where $\lambda_i > 0$ are the eigenvalues of the Hessian (analogous to spring stiffness $k$). 
To ensure the fastest convergence without oscillation (Critical Damping), we must tune the hyperparameters $m$ and $b$ such that $b^2 = 4m\lambda_i$. Since $\lambda_i$ varies across dimensions, perfect critical damping in all directions is impossible, leading to the rich field of adaptive momentum algorithms (like Adam) in modern AI research.

\newpage
\section{Part 7: Common Variations and Untested Concepts}

\subsection{Missing Concept: Forced Vibrations and Resonance}
The provided homework strictly covered \textit{free} vibrations (where the right side of the ODE is 0). However, Section 4.10 introduces \textbf{Forced Vibrations}, where an external force $F_{\text{ext}}(t) = F_0 \cos(\gamma t)$ is applied. The ODE becomes nonhomogeneous:
\[ my'' + by' + ky = F_0 \cos(\gamma t) \]
A critical phenomenon occurs when the system is undamped ($b=0$) and the forcing frequency $\gamma$ exactly matches the natural frequency $\omega = \sqrt{k/m}$. This is called \textbf{Pure Resonance}.

\textbf{Example Problem:} Find the general solution to $y'' + 9y = 5 \cos(3t)$.

\textbf{Solution:}
1. \textbf{Homogeneous Solution:} $r^2 + 9 = 0 \implies r = \pm 3i$.
\[ y_h(t) = C_1 \cos(3t) + C_2 \sin(3t) \]
Notice the natural frequency is $\omega = 3$.

2. \textbf{Particular Solution (Method of Undetermined Coefficients):}
Since the forcing function $5 \cos(3t)$ is already in the homogeneous solution, we must multiply our guess by $t$ to ensure linear independence.
\[ y_p(t) = t(A \cos(3t) + B \sin(3t)) \]
Differentiating:
\[ y_p'(t) = A \cos(3t) + B \sin(3t) + t(-3A \sin(3t) + 3B \cos(3t)) \]
\[ y_p''(t) = -6A \sin(3t) + 6B \cos(3t) - 9t(A \cos(3t) + B \sin(3t)) \]
Substitute into $y'' + 9y = 5 \cos(3t)$:
\[ -6A \sin(3t) + 6B \cos(3t) - 9y_p + 9y_p = 5 \cos(3t) \]
\[ -6A \sin(3t) + 6B \cos(3t) = 5 \cos(3t) \]
Equating coefficients: $-6A = 0 \implies A = 0$, and $6B = 5 \implies B = 5/6$.
\[ y_p(t) = \frac{5}{6}t \sin(3t) \]

3. \textbf{General Solution:}
\[ y(t) = C_1 \cos(3t) + C_2 \sin(3t) + \frac{5}{6}t \sin(3t) \]
\textit{Notice the $t$ in the particular solution. As $t \to \infty$, the amplitude of the oscillation grows to infinity. This mathematical result explains why bridges can collapse if soldiers march in step at the bridge's natural frequency!}

\newpage
\section{Part 8: Advanced Diagnostic Testing: "Find the Flaw"}

Below are 5 problems with solutions that contain a single, subtle, yet critical flaw. Your task as the grader is to find the error, explain why it is wrong in one sentence, and provide the correct step and final answer.

\subsection{Flawed Problem 1}
\textbf{Problem:} An undamped spring-mass system has the equation of motion $y(t) = -3 \cos(2t) + 4 \sin(2t)$. Write this in the form $y(t) = A \sin(2t + \phi)$.
\textbf{Flawed Solution:}
1. Amplitude $A = \sqrt{(-3)^2 + 4^2} = \sqrt{9 + 16} = 5$.
2. Phase angle $\tan \phi = \frac{C_1}{C_2} = \frac{-3}{4} = -0.75$.
3. $\phi = \arctan(-0.75) \approx -0.6435$ radians.
4. Final Answer: $y(t) = 5 \sin(2t - 0.6435)$.

\subsection{Flawed Problem 2}
\textbf{Problem:} Find the general solution to the critically damped system $y'' + 12y' + 36y = 0$.
\textbf{Flawed Solution:}
1. Characteristic equation: $r^2 + 12r + 36 = 0$.
2. Factoring: $(r + 6)^2 = 0 \implies r = -6$.
3. Since it is a repeated root, the general solution is a linear combination of the root.
4. Final Answer: $y(t) = C_1 e^{-6t} + C_2 e^{-6t}$.

\subsection{Flawed Problem 3}
\textbf{Problem:} A damped system has the equation of motion $y(t) = e^{-2t} \sin(3t)$. Find the time $t > 0$ when the mass first reaches its maximum displacement.
\textbf{Flawed Solution:}
1. To find the maximum displacement, set $y(t) = 0$.
2. $e^{-2t} \sin(3t) = 0$.
3. Since $e^{-2t}$ is never zero, we must have $\sin(3t) = 0$.
4. $3t = \pi \implies t = \pi/3$.
5. Final Answer: The maximum displacement occurs at $t = \pi/3$ seconds.

\subsection{Flawed Problem 4}
\textbf{Problem:} Determine if the system $2y'' + 5y' + 4y = 0$ is underdamped, critically damped, or overdamped.
\textbf{Flawed Solution:}
1. Identify the constants: $m = 2$, $b = 5$, $k = 4$.
2. Calculate the discriminant: $b^2 + 4mk$.
3. $5^2 + 4(2)(4) = 25 + 32 = 57$.
4. Since $57 > 0$, the system has two distinct real roots.
5. Final Answer: The system is overdamped.

\subsection{Flawed Problem 5}
\textbf{Problem:} Find the derivative $y'(t)$ for the underdamped motion $y(t) = 5e^{-3t} \cos(4t)$.
\textbf{Flawed Solution:}
1. Apply the derivative rules to the function.
2. The derivative of $5e^{-3t}$ is $-15e^{-3t}$.
3. The derivative of $\cos(4t)$ is $-4\sin(4t)$.
4. Multiply them together.
5. Final Answer: $y'(t) = 60e^{-3t} \sin(4t)$.

\end{document}